\section{Analysis}

\subsection{Kernel Approximation Quality}

\textbf{Figure 8.1} compares the empirical kernel $\phi(q)^\top \phi(k)$ 
against the true softmax kernel $\exp(q^\top k / \sqrt{d_h})$ and the 
Gaussian kernel $\exp(-\|q - k\|^2 / 2)$.

\textbf{Observation.} The RFF approximation is closer to the Gaussian kernel 
than the softmax kernel, as expected from theory.

\subsection{Effect of Feature Dimension}

\textbf{Figure 8.2} shows accuracy vs $F$. Approximation error scales as 
$O(1/\sqrt{F})$, and accuracy improves with larger $F$ up to a point.

\begin{table}[H]
\centering
\small
\begin{tabular}{lll}
\toprule
\textbf{$F$} & \textbf{Approx. err} & \textbf{Accuracy} \\
\midrule
64 & ~0.15 & {{D\_F64}} \\
128 & ~0.10 & {{D\_F128}} \\
256 & ~0.07 & {{D\_ACC}} \\
512 & ~0.05 & {{D\_F512}} \\
\bottomrule
\end{tabular}
\end{table}

\subsection{Preservation Gate Analysis}

\textbf{Figure 8.3} shows the learned $\alpha$ per layer for DistilBERT.

\begin{table}[H]
\centering
\small
\begin{tabular}{ll}
\toprule
\textbf{Layer} & \textbf{$\alpha$} \\
\midrule
0 & {{D\_L0\_ALPHA}} \\
1 & {{D\_L1\_ALPHA}} \\
... & ... \\
5 & {{D\_L5\_ALPHA}} \\
\bottomrule
\end{tabular}
\end{table}

\textbf{Observation.} Early layers prefer attention (high $\alpha$); later layers 
{{LATE\_LAYER\_OBS}}.

\subsection{Why QR Helps}

Sampling $\omega_i \sim \mathcal{N}(0, I)$ gives an unbiased estimator 
with variance $\sigma^2 = O(1/F)$. Orthogonalizing reduces this to 
$O(1/F^2)$, giving faster convergence.

\textbf{Table 8.1: Empirical kernel error.}

\begin{table}[H]
\centering
\small
\begin{tabular}{llllll}
\toprule
\textbf{Init} & \textbf{} & \textbf{} & \textbf{(q)  (k) -} & \textbf{} & \textbf{$_2$} \\
\midrule
Random & {{ERR\_RANDOM}} &  &  &  &  \\
QR & {{ERR\_QR}} &  &  &  &  \\
\bottomrule
\end{tabular}
\end{table}

\subsection{LayerNorm Before Projection}

Without LN, the input to the random projection has varying scale, which 
biases the kernel approximation. LN normalizes this.

\textbf{Figure 8.4} shows the distribution of $f = \tilde{x} W$ with and 
without LN.

\subsection{Causal Cumsum Memory}

For decoders, the cumsum over the outer product materializes a large 
tensor. \textbf{Table 8.2} shows peak memory:

\begin{table}[H]
\centering
\small
\begin{tabular}{ll}
\toprule
\textbf{L} & \textbf{Peak memory} \\
\midrule
128 & {{D\_MEM128}} MB \\
256 & {{D\_MEM256}} MB \\
512 & {{D\_MEM512}} MB \\
\bottomrule
\end{tabular}
\end{table}

For $L > 512$, memory becomes a bottleneck. See Appendix D for 
optimization.

\subsection{Comparison to Other Linear Attention}

\textbf{Table 8.3: LEKAv1 vs other linear attention.}

\begin{table}[H]
\centering
\small
\begin{tabular}{lll}
\toprule
\textbf{Method} & \textbf{DistilBERT} & \textbf{Speed @ L=1024} \\
\midrule
full & 90.71\% & 1$\times$ \\
Linear Transformer & {{LT\_D}} & {{LT\_SP}}$\times$ \\
Performer & {{PERF\_D}} & {{PERF\_SP}}$\times$ \\
\textbf{LEKAv1} & \textbf{{{D\_ACC}}} & \textbf{{{D\_SP1024}}$\times$} \\
\bottomrule
\end{tabular}
\end{table}

\subsection{Failure Modes}

\subsubsection{Extractor Failure}

The current \texttt{copy\_weights} raises \texttt{RuntimeError} if UniversalAdapter fails. 
This \textbf{kills the module} and can lead to collapse.

\subsubsection{Cumsum OOM}

At large L, the cumsum materializes a huge tensor and can cause OOM.

\subsubsection{Preservation Gate Saturation}

If $\alpha$ saturates to 0, the attention branch is ignored. If $\alpha$ saturates to 1, 
the residual is ignored.

\subsection{Summary of Insights}

\begin{enumerate}
\item \textbf{RFF works} for attention approximation.
\item \textbf{QR projection reduces variance.}
\item \textbf{Head-wise LN stabilizes input.}
\item \textbf{Preservation gate provides safety.}
\item \textbf{Linear scaling enables long-context inference.}

\end{enumerate}
